% commentaire.tex — f163v, vv. 80–122 (conversion initiale des travaux antérieurs)
\comm{80}{ἐπεὶ τρίτη 'στὶ Θεσμοφορίων ἡ μέση.}{Le manuscrit de Ravenne donne au vers 80 : « puisque c'est le troisième [jour] des Thesmophories, celui du milieu. » Problème : si les Thesmophories duraient bien trois jours, celui du milieu ne peut en aucune façon être le troisième. Cette contradiction eut la fâcheuse conséquence de déchaîner des torrents de commentaires (à commencer par celui du scholiaste) ; on a inventé, en contradiction avec tous les témoignages\nc{Outre le scholiaste, qui \italique{sait} que les Thesmophories ne duraient que trois jours – sinon il ne prendrait pas la peine de tant ergoter -- et la glose dubitative d'Hésychius (τ 1440), citons : Alciphron 2,37.2 ; 3,39 ; et la « \anglais{speculative restauration} » du papyrus n° 3540.1-2, \italique{The Oxyrhynchus Papyri}, L (1983), ed. E. W. Handley}, des Thesmophories de quatre jours à grand renfort d'érudition et de savantes notes, et ce sous le prétexte qu'il ait pu y avoir avant ces trois jours d'autres festivités ailleurs (surprenant, non ? mais il suffit de prendre un calendrier des fêtes grecques pour constater qu'avec ce genre de raisonnements, on pourrait tout aussi bien \italique{prouver} que les Thesmophories duraient 365 jours par an ; au reste, il faudrait encore montrer que le milieu d'un ensemble de quatre éléments en est le \nieme{3} ; pour finir – \italique{last but not least} –, faire préciser à Aristophane quel jour tombe ἡ μέση serait un détail inepte tant cela était connu de tous les Athéniens). Le texte avait été pourtant brillamment rétabli par August Nauck\nc{in \italique{Mélanges gréco-romains}, St.-Pétersburg (1866), t. II, pp. 468-71 ; voir aussi pp. 344-5 pour des témoignages sur les \italique{Thesmophories}.} et intégré avec explications dès l'édition aux Belles Lettres de Victor Coulon ; cela n'empêche pas la faute d'apparaître dans nombre d'éditions modernes. Venons-en donc à la bonne correction : l'apparition du τρίτη, \italique{troisième}, est due à la mauvaise lecture d'un Π' (abréviation de περί ou περ) en Γ' (\italique{troisième} lettre de l'alphabet) dans un ancien papyrus\nc{La corruption doit être ancienne puisqu'elle – conjointe au bavardage du scholiaste – rendait déjà perplexe Hésychius (τ 1440): τρίτη Θεσμοφορίων· ζητεῖται πῶς ἅμα μὲν λέγει τρίτην Θεσμοφορίων εἶναι, ἅμα δὲ μέσην, τεττάρων οὐσῶν ἡμερῶν.} – cette abréviation est fréquente, et source d'erreurs aussi dans l'autre sens ; ainsi a-t-il fallu en maintes occasions rétablir un `τρίτη' à partir d'un gamma\nc{Pour rester avec notre cher Aristophane, un `τρεῖς' a dû être restitué au vers 1507 des \italique{Guêpes} là où tous les manuscrits comportent un `γε'.}. Bref, voici l'histoire de notre début de \nieme{80} vers, très ingénieusement reconstituée : $$\overset{\text{\footnotesize{\allemand{Urtext}}}}{\text{ἐπείπερ ἐστὶ}}\xrightarrow{\text{onciales}}\text{ΕΠΕΙΠ'ΕCΤΙ}\xrightarrow{\text{corruption}}\text{ΕΠΕΙΓ'ΕCΤΙ}\xrightarrow{\text{interprétation}}\overset{\text{\footnotesize{\italique{Ravennas}}}}{\text{ἐπεὶ τρίτη ᾽στὶ.}}$$}
\comm{83}{κἀν Θεσμοφόροιν.}{Le scholiaste précise que c'est une Ἀττικὴ σύνταξις καὶ δυϊκὴ, ἀντὶ τοῦ `ἐν τῷ τῶν θεσμοφόρων ἱερῷ' ; les \italique{deux} déesses thesmophores étant Déméter et sa fille Perséphone.}
\comm{86–7}{νὴ τὸν Ποσειδῶ καὶ δί᾽ και᾽ ἂν πάθοις / ἀτὰρ τίν᾽ ἐκ ταύτης σὺ μηχανὴν ἔχεις;}{Ne voyant pas à quel nom féminin renvoyait `ἐκ ταύτης', on a retouché le texte du vers 87 à l'envi (ἐκ τούτων Bergler, εἰς ταύτας Sandbach, \&c.). 
C'est à mon avis aller un peu vite en besogne alors qu'il y a justement un problème textuel au vers précédent : il manque une syllabe, probablement au niveau du fautif `δί᾽ και᾽'. Je fais l'hypothèse que c'est justement à cet endroit que se trouvait ce nom féminin, repris par le `ἐκ ταύτης' du vers suivant. De quoi Euripide peut-il se mettre \italique{hors} en déployant quelque machine ? Réponse : tel le Strepsiade des \italique{Nuées}, c'est à la justice qu'il veut échapper. Au vu du texte conservé, la correction est évidente : il faut substituer au `δί᾽ και᾽' un `δικαίαν' ; la chute de l'`αν' s'expliquant par la courante erreur de «~saut du même au même~» (`δικαίαν' étant directement suivi de la particule `ἂν'). 
Lisons donc :
\citComm{νὴ τὸν Ποσειδῶ καὶ δικαίαν ἂν πάθοις.\\
ἀτὰρ τίν᾽ ἐκ ταύτης σὺ μηχανὴν ἔχεις;\\
\francais{Sainte Mer ! ce serait en justice que tu pâtirais.\\
Mais pour y échapper, de quelle machine disposes-tu ?}}
Cette correction est d'autant plus intéressante que ἡ δικαία, mot rare, est un synonyme poétique de δίκη\nc{Le mot est attesté dans l'\EtymMagnum à l'entrée Ἀθηνᾶ : le glosateur y recense Ἀθηναία, qui est à Ἀθήνη ce que δικαία est à δίκη (Ἀθήνη Ἀθηναία, ὡς δίκη δικαία).} ; Mnésiloque retourne-t-il contre Euripide une de ses propres sentences\nc{Tout au long de la pièce, Euripide est pris au piège de ses propres γνῶμαι (aux vers 194, 275-276, etc). Est-ce aussi le cas ici ? Quoi qu'il en soit, on trouve des expressions très proches aux vers 1379 des \italique{Nuées} : \citComm{Φειδιππίδης -- νὴ τὸν Δί᾽ ἐν δίκῃ γ᾽ ἄν,}ou 403 de \italique{Lysistrata} :\citComm{Πρόβουλος -- νὴ τὸν Ποσειδῶ τὸν ἁλυκὸν δίκαιά [δικαία ?] γε.}} ?\\
Wilamowitz\nc{Au commentaire du vers 681 de \italique{La Folie d'Héraclès} : second volume de son \italique{Euripides Herakles erklärt von Ulrich von Wilamowitz-Moellendorff}. Wilamowitz cite \italique{en passant} le vers 86 des \italique{Thesmophories} ainsi corrigé sans même préciser la correction -- dont il semble bien être l'auteur -- tant elle devait lui sembler aller de soi.} avait déjà proposé cette émendation sans qu'elle soit jamais reprise dans les éditions publiées depuis -- à une notable exception : celle de Prato.}
\comm{87}{μηχανὴν.}{Au sens imagé, une μηχανή est un \italique{expédient}, un \italique{tour}, un \italique{stratagème}. Certes, Euripide est un roué, un filou, un retors... mais Euripide est avant tout un homme de théâtre ! Il dispose de \italique{machineries} dont il use et abuse\nc{Euripide recourut à la μηχανή massivement dans ses œuvres tardives. Aristote, dans sa \italique{Poétique} (XV, 1454a37), le lui reproche pour sa \italique{Médée} : Euripide privilégia un dénouement ἀπὸ μηχανῆς plutôt que par l'histoire elle-même : φανερὸν οὖν ὅτι καὶ τὰς λύσεις τῶν μύθων ἐξ αὐτοῦ δεῖ τοῦ μύθου συμβαίνειν, καὶ μὴ ὥσπερ ἐν τῇ Μηδείᾳ ἀπὸ μηχανῆς. Aristophane parodie en maints endroits cet artifice de théâtre. Pour un dossier plus complet, aller voir l'article «~Machines et \italique{deus ex machina} dans les spectacles comiques~» de Leonardo Fiorentini, in \italique{L'appareil scénique dans les spectacles de l'Antiquité} (2013).}. Ce sont ces deux sens de μηχανή, le propre et le figuré, qui étaient instantanément saisis du spectateur du \nieme{V} siècle av. J.-C. Conformité du langage français avec le grec : le mot « truc », qui désignait au \nieme{XVII} siècle une machine de théâtre, s'est approprié le même sens imagé.}
\comm{90–1}{ἐκκλησιάσοντ᾽ ἐν ταῖς γυναιξὶν κἂν δέῃ / λέξονθ᾽ ὑπὲρ ἐμοῦ.}{Je reprends la correction de Markland : substituer χἂν ($=$ καὶ ἃ ἂν) à κἂν ($=$ καὶ ἐὰν) qui présente un sens bien meilleur : Euripide ne dit certes pas « qu'au besoin » ou « s'il le faut » (κἂν) Agathon interviendra mais bien plutôt qu'il lui faut intervenir autant qu'il le pourra (χἂν). Euripide a sa vie en jeu – si l'on peut s'exprimer ainsi – et sait bien le courroux de la gent féminine ; l'inaction aurait forcément une issue fatale... heureuse pour Athènes et les Athéniennes mais un peu plus fâcheuse pour sa personne. 
Précisons que ἐν a ici le sens technique que donne le Bailly (A, III.10) : « lorsqu'il s'agit d'un orateur, d'un accusé, ἐν correspond au français \italique{devant} » (cf. les \italique{Acharniens} v. 630 ; les \italique{Guêpes} v. 1185 ; le fragment 161 K.-A.) ; il faut donner un sens actif au `ἐκκλησιάσοντ᾽ ἐν ταῖς γυναιξὶν' : « pour intervenir à l'assemblée devant les femmes ».}
\comm{94}{ἡμέτερος ὁ πυραμοῦς.}{L'expression proverbiale signifie : « à nous la palme de la victoire ». En effet, on offrait un gâteau de miel et de froment (πυραμοῦς), à titre de récompense, à celui qui avait le mieux résisté au sommeil dans une παννυχίς (scholiaste des \italique{Cavaliers}, au v. 277 où l'expression apparaît à l'identique ; Athénée 647b-c, 668c). (Note de Taillardat, §847.)}
\comm{96}{οὑκκυκλούμενος}{($=$ ὁ ἐκκυκλούμενος). Agathon apparaît sur l'eccyclème, « machine de théâtre munie de roues qui, en tournant, faisait voir aux spectateurs ce qui était censé se passer dans une maison. » (Chantraine.) 
Le scholiaste l'avait noté :\citComm{ἐπὶ ἐκκυκλήματος γὰρ φαίνεται.}}
\comm{98}{Κυρήνην δ᾽ ὁρῶ.}{Cyrène était une célèbre courtisane évoquée aussi aux vers 1327-1328 des \italique{Grenouilles} : 
\citComm{ἀνὰ τὸ δωδεκαμήχανον / Κυρήνης μελοποιῶν\\
\francais{$[$Euripide$]$ compose ses chants à la mode de Cyrène aux douze postures.}}Le scholiaste y est plus loquace que le nôtre et nous apprend le trait distinctif de cette Cyrène :\citComm{Κυρήνη τις ἑταίρα ἐπίσημος, δωδεκαμήχανος ἐπικαλουμένη διὰ τὸ τοσαῦτα σχήματα ἐν τῇ συνουσίᾳ ποιεῖν.}Taillardat (§734) rapproche pertinemment ce sobriquet du titre d'un ``traité'' de Paxamos : Δωδεκάτεχνον, à propos duquel la \italique{Souda} (à l'entrée Πάξαμος) dit : ἔστι δὲ περὶ αἰσχρῶν σχημάτων. Que moque donc Aristophane en rapprochant, non cette fois la muse du poète comme dans les \italique{Grenouilles}, mais Agathon lui-même de Cyrène ? Bien-sûr son air efféminé ! bien-sûr son travestissement !.. mais pourquoi Cyrène en particulier ? Réponse : pour ce qu'elle a de remarquable (ἐπίσημος) qui lui valut son surnom ; il ne me paraît pas douteux qu'Agathon, qui vient de se faire rouler par l'eccyclème, apparaît affalé sur son lit dans une position suggestive, ou du moins qui peut prêter à la confusion de Mnésiloque. (Cf. la raillerie -- μῶν βινεῖσθαι; -- au vers 50.)}
\comm{99}{μινυρισμός.}{Une scholie du \italique{Ravennas} semble avoir gardé la trace d'une ancienne didascalie : μινυρισμός était indiqué marginalement :
\citComm{μεταξὺ τῶν δυοῖν ἀξιοῦσί τινες γράφειν `μινυρισμός', ὡς πολλὰ τοιαῦτα παρεπιγράφεται.}}
\comm{100}{μύρμηκος ἀτραπούς.}{Willems, en accord avec le scholiaste, se propose d'expliquer ainsi l'expression μύρμηκος ἀτραπούς : « par allusion à la marche saccadée et tortueuse des fourmis. Image familière aux Anciens. Dans l'\italique{Anthologie}, XI, 78, il est question des galeries obliques et droites, λόξα καὶ ὀρθά, d'une fourmilière. Voir aussi Plutarque, \italique{de Solertia animal.}, II, p. 968a. Plaute a dit : \latin{formicinum gradum movere}. » Taillardat (§784) abonde en son sens avec d'autres exemples à l'appui. Cependant, pour suggérer l'image de \italique{sentiers de fourmis}, n'aurait-il pas fallu, en français comme en grec, un pluriel (μυρμήκ\textbf{ων} ἀτραπούς) ?\\
Hésychius ne le comprenait pas autrement et voyait dans `μύρμηκος' un nom propre (μ 1904) :
\citComm{Μύρμηκος ἀτραπούς· Ἀθήνησιν ἐν Σκαμβωνιδῶν ἔστι Μύρμηκος ἀτραπός, ἀπὸ ἥρωος Μύρμηκος ὀνομαζομένη.\\
\francais{Dans [le dème] athénien des Scambônides [σκαμβός signifie \italique{courbe}, \italique{tortu}] se trouve un sentier de Myrmex, qui tire son nom du héros Myrmex.}}
Que l'image soit tirée de la fourmilière ou d'un nom propre, la critique n'en reste pas moins d'ordre musical\nc{L'image des lignes brisées de la fourmilière, μυρμηκιά, servit en musique pour décrire des trilles ou des arpèges. On trouve l'expression `ἐκτραπέλους μυρμηκιάς' chez le Pseudo-Plutarque, \italique{De Musica}, 30 ; il s'agit du fragment 155(v.23) de Phérécrate dans l'édition Kassel-Austin. À noter que le même Phérécrate a composé une comédie intitulée les μυρμηκάνθρωποι.}.
Philoxène de Cythère, un représentant de cette «~nouvelle musique~» dénoncée par Aristophane, reçut d'ailleurs le surnom de Μύρμηξ (voir \italique{Souda} φ 393, à l'entrée Φιλόξενος). Le persiflage de Mnésiloque ne serait-il pas une allusion audit Philoxène, alias Myrmex ?}
\comm{101}{Χθονίαιν}{est un génitif se rapportant à ἱεράν. Cf. \italique{Ploutos}, 937, et \italique{Caval.}, 301. (Note de Willems.)}
\comm{101–29}{Agathon chante un fragment de tragédie.}{Sujet : un coryphée et un chœur de jeunes filles troyennes célèbrent en strophes alternées Léto et ses deux rejetons. Il est très évident qu'Agathon chante seul. Les arguments fournis par M. P. Mazon\nc{in \italique{Essai sur la composition des comédies d'Aristophane} (1904), p. 127-128.} sont à cet égard sans réplique. (Note de Willems.)\\
Suidas (χ 403) nous confirme qu'Agathon avait l'habitude de chanter de la sorte : 
\citComm{Χορικός· ὁ τοῦ χοροῦ ἐξάρχων. τοιοῦτος ἦν Ἀγάθων, ὃς ὑποκριτικὰ μέλη ἐποίει καὶ αὐτὸς ὑποκρίνεται\\
\francais{Chorique : celui qui mène le chœur. Tel Agathon, qui avait l'habitude de créer des chants dramatiques et de les interpréter lui-même.}}L'on pas manqué de remarquer en quoi les ruptures de rythmes (παράρυθμ' εὔρυθμα) ou le contexte ionique de cette parodie d'Agathon en font un morceau efféminé et suave, ce qui est tout-à-fait juste quoiqu'il nous manque la musique pour juger sur pièces\nc{À ce propos, cf. les \italique{Grenouilles} v. 1299 où Eschyle dit s'être écarté du pré des Muses de Phrynichos – Muses par trop efféminées et ionisantes ; de même, l'\italique{Assemblée des femmes} (v. 883).} ; mais beaucoup trop de commentateurs se sont arrêté à ces considérations\nc{Le style d'Agathon est certes efféminé mais beaucoup moins caricatural qu'on a voulu le faire croire. Cf. le fragment 178 (Γηρυτάδης) de l'édition Kassel-Austin : Aristophane y emploie l'adjectif Ἀγαθώνειος dans un sens loin d'être péjoratif d'après les commentaires anciens (Zénobius I, 2 ; Hésych. α 281 ; Phot. α 83 ; Suid. α 125 ; \&c.) :\citComm{Ἀγαθώνιος αὔλησις: ἡ μαλακὴ καὶ ἐκλελυμένη: ἢ ἡ μήτε χαλαρὰ, μήτε πικρὰ, ἀλλ' εὔκρατος καὶ ἡδίστη.\\
Un air de flûte agathonien : doux et détendu ; ni mol ni strident, mais tempéré et très doux.}}. 
(Opinion au demeurant est bien pratique puisqu'elle dispense d'avoir à porter trop d'attention sur le texte.) Ce fut pas le cas du regretté Alain Blanchard avait excellemment analysé la parodie d'Agathon dans son article : \italique{La parodie d'Agathon et la structure des }``Bucoliques''\italique{ de Virgile}\nc{\italique{Dans l'ouvroir du poète}, \italique{Hellenica} (2008), pp. 35-45. J'eus la chance de recevoir d'Alain Blanchard peu de temps avant sa mort de vifs encouragements, source d'énergie nécessaire à l'abattage de la somme de travail que représente ce commentaire des \italique{Thesmophories}.}. Blanchard décela dans ce chant une structure concentrique, en cela analogue aux \italique{Bucoliques} de Virgile ou encore aux cantates de Bach ; le point important étant que « la reconnaissance de ce schéma concentrique » est d'une grande « valeur exégétique ». Laissons-lui la parole : « la coryphée invite le chœur à célébrer Apollon, Artémis et leur mère Léto, et le chœur, docilement, obéit à cette injonction. En réalité, l'étude de composition, une fois admise l'existence d'une structure concentrique conduit à plus de subtilité.
Le parallélisme entre le premier élément (vv. 101-103) et le neuvième (vv. 126-129) réside dans le thème de la lumière (lumière des torches portées par le chœur, lumière qui brille en réponse dans les yeux divins), thème qui précède un appel à la célébration d'un dieu – l'appel étant plus solennel, plus brusque et plutôt à titre de rappel dans le neuvième. Le parallélisme entre le deuxième élément (104-106) et le huitième (123-125) réside dans le thème du σέβεσθαι, introduit en mineur à la fin du deuxième élément et qui éclate sur le mode majeur au début du huitième.
Le parallélisme entre le troisième élément (107-110) et le septième (120-122) réside dans le caractère national des divinités honorées : caractère national sensible surtout dans le troisième élément (Apollon constructeur des remparts de Troie), mais présent, plus discrètement, dans le septième avec les Charites phrygiennes.
Enfin, le thème de la majesté divine (majesté d'Apollon comme dispensateur de la gloire) entre le quatrième élément (111-113) et le sixième (117-119). Le cinquième (114-116), central, est un appel à honorer Artémis, divinité féminine. Dans la première partie de la parodie, les trois interventions du coryphée concernent des divinités armées de l'arc qu'il faut honorer avec un esprit libre et d'une voix forte (103), mais le coryphée s'adresse à un chœur de jeunes filles qui, pour ainsi dire, retourne complètement la nature de ces divinités dont on aperçoit alors le caractère ambigu. Déjà, aux appels à un culte viril qu'on leur faisait, elles avaient répondu, avec une soumission apparente bien féminine, par une profession de piété plus neutre (σεβίσαι, 106). Puis, avec elles, Apollon cesse d'être le dieu archer pour être le coryphée des Muses (112). Qu'on leur propose d'honorer Artémis, il n'est plus question avec elle de chasse – occupation virile – mais de maternité : maternité glorieuse de Léto (118), maternité refusée par elle-même par la déesse Vierge (119). Le ton est donné pour la seconde partie du poème où le chœur, à l'appel du coryphée, honorera, après Artémis, sa mère Léto et la cithare. Dans un cadre aussi favorable, le chœur peut même se permettre une concession – nous parlerions plutôt d'un appel spontané et aventureux du cœur – avec l'expression ἄρσενι βοᾷ, \italique{mâle accent}, en 125 qui rappelle en la précisant l'invitation faite en 103b par le coryphée. Mais qu'on invite de nouveau à honorer une divinité masculine comme Apollon : c'est la maternité de Léto qui est chantée (130). »\\
Fera-t-on la même erreur à propos de ce chant que les premiers auditeurs du sonnet d'Oronte en prétendant : « Il est bon » ? C'est en tout cas ce qu'il nous semble à l'analyse de ce petit bijou de malice au subtil renversement : d'un banal hommage au fondateur d'une cité, l'on passe doucement et insidieusement à un morceau d'érotisme sur fond de maternité frustrée, le tout échafaudé sur une structure solidement ancrée\nc{Cf. pour plus de détails l'article d'A. Blanchard avec les chiffres à la clef.}. Rien d'étonnant donc à ce que Mnésiloque en soit tout émoustillé ; rien d'étonnant non plus à ce que, par la bouche de Dionysos, Aristophane rende hommage à Agathon après sa mort : ἀγαθὸς ποιητὴς καὶ ποθεινὸς τοῖς φίλοις, « un bon poète apprécié de ses amis » (\italique{Grenouilles}, v. 84). Pour finir, l'élaboration complexe du morceau ainsi mise au jour par Blanchard fait voler en éclat la thèse fort répandue d'une improvisation de la part d'Agathon ; d'ailleurs l'on sait par Aristote qu'Agathon avait eu l'ambition dans une de ses tétralogies de traiter intégralement du sac de Troie, ambition qui l'avait mené à l'échec : (\italique{Poétique}, XVIII, 56a16-19) :
\citComm{ὅσοι πέρσιν Ἰλίου ὅλην ἐποίησαν καὶ μὴ κατὰ μέρος ὥσπερ Εὐριπίδης, <ἢ> Νιόβην καὶ μὴ ὥσπερ Αἰσχύλος, ἢ ἐκπίπτουσιν ἢ κακῶς ἀγωνίζονται, ἐπεὶ καὶ Ἀγάθων ἐξέπεσεν ἐν τούτῳ μόνῳ.\\
\francais{Tous les poètes qui veulent traiter la ruine de Troie entièrement et non pas par parties comme Euripide, ou la légende de Niobé tout entière au lieu de faire comme Eschyle, tantôt échouent complètement, tantôt font mauvaise figure dans les concours. Ainsi, Agathon échoua sur ce seul point.}}Je me plais à imaginer Aristophane, révolté devant cet échec injuste, remettre sur scène ce qu'il considérait comme la perle (érotique) du cycle troyen d'Agathon (c'est en tout cas ce sens de la justice qui l'avait poussé à vouloir faire rejouer ses \italique{Nuées}, mésestimées par les mauvais jurés athéniens).}
